\documentclass[12pt, a4paper]{report}
\usepackage[utf-8]{inputenc}
\usepackage[spanish]{babel}
\usepackage{graphicx}
\usepackage{minted}
\usepackage[hidelinks]{hyperref}
\usepackage{fancyhdr}
\usepackage{geometry}
\usepackage{xcolor}
\usepackage{listings}
\usepackage{amsmath}
\usepackage{amssymb}
\usepackage{setspace}
\usepackage{titlesec}

% Configuración de márgenes
\geometry{
    left=2.5cm,
    right=2.5cm,
    top=2.5cm,
    bottom=2.5cm
}

% Configuración de fuentes y espaciado
\onehalfspacing

% Configuración de headers
\pagestyle{fancy}
\fancyhf{}
\fancyhead[R]{\thepage}
\fancyhead[L]{Análisis de Patrones de Diseño - WallpaWawqi}
\renewcommand{\headrulewidth}{0.5pt}

% Configuración de títulos
\titleformat{\chapter}[display]
    {\Huge\bfseries}
    {Capítulo \thechapter}{20pt}{\Huge}
\titleformat{\section}
    {\Large\bfseries}
    {}{0pt}{}
\titleformat{\subsection}
    {\large\bfseries}
    {}{0pt}{}

% Configuración de minted
\setminted{
    linenos=true,
    breaklines=true,
    fontsize=\small,
    bgcolor=lightgray!20
}

\usemintedstyle{monokai}

\title{
    \vspace*{\fill}
    {\Huge \textbf{Análisis Detallado de Patrones de Diseño}}\\
    \vspace{1cm}
    {\LARGE Proyecto: WallpaWawqi}\\
    \vspace{0.5cm}
    {\large Backend - Sistema de Gestión de Restaurante}
    \vspace*{\fill}
}

\author{
    \Large Análisis de Arquitectura y Diseño\\
    \large Sistema de Información WallpaWawqi
}

\date{\today}

\begin{document}

% Portada
\maketitle

% Página en blanco
\newpage
\thispagestyle{empty}
\mbox{}

% Tabla de contenidos
\newpage
\tableofcontents

% Introdución
\chapter{Introducción}

El presente informe realiza un análisis detallado de los patrones de diseño implementados en el backend del proyecto WallpaWawqi, un sistema de gestión integral para restaurantes desarrollado en Java.

\section{Contexto del Proyecto}

WallpaWawqi es una aplicación que gestiona la operación de un restaurante, incluyendo:
\begin{itemize}
    \item Gestión de empleados por rol
    \item Catálogo de platillos y menú
    \item Persistencia de datos en formato JSON
    \item API REST para consultas y registro de platillos
\end{itemize}

\section{Objetivo}

Identificar, documentar y analizar los patrones de diseño implementados en la lógica del backend, evaluando su efectividad, coherencia arquitectónica y aplicabilidad a los requisitos del sistema.

\section{Alcance}

El análisis se enfoca exclusivamente en la lógica backend, excluyendo la interfaz web (carpeta webAPP). Se examinan:
\begin{itemize}
    \item Estructura de clases y herencia
    \item Servicios de persistencia
    \item Manejo de solicitudes HTTP
    \item Uso de genéricos y programación funcional
\end{itemize}

\newpage

\chapter{Estructura del Proyecto}

\section{Arquitectura General}

La arquitectura del backend se organiza en las siguientes capas:

\begin{itemize}
    \item \textbf{Capa de Modelos}: Clases que representan entidades de dominio
    \item \textbf{Capa de Servicios}: Lógica de persistencia y acceso a datos
    \item \textbf{Capa de Presentación}: Servidor HTTP que expone endpoints REST
\end{itemize}

\section{Componentes Principales}

\subsection{Jerarquía de Empleados}

\begin{verbatim}
Empleado (abstract)
  ├── Cocinero
  ├── Mozo
  └── PersonalDeLimpieza
\end{verbatim}

\subsection{Entidades de Negocio}

\begin{itemize}
    \item \textbf{Platillo}: Representa los platos disponibles en el menú
    \item \textbf{Empleado}: Entidad base para todos los tipos de empleados
\end{itemize}

\subsection{Servicios}

\begin{itemize}
    \item \textbf{JsonPersistenceService}: Gestiona la persistencia genérica de datos
\end{itemize}

\newpage

\chapter{Patrones de Diseño Identificados}

\section{1. Template Method Pattern}

\subsection{Descripción}

El patrón Template Method define el esqueleto de un algoritmo en una clase base, permitiendo que las subclases sobrescriban pasos específicos sin cambiar la estructura del algoritmo.

\subsection{Aplicación en WallpaWawqi}

La clase abstracta \texttt{Empleado} implementa este patrón al:

\begin{itemize}
    \item Definir una estructura común para todos los empleados (propiedades básicas)
    \item Proporcionar métodos plantilla como \texttt{registrarAsistencia()}
    \item Permitir que cada subclase especializada implemente sus propios comportamientos específicos
\end{itemize}

\subsection{Diagrama Conceptual}

\begin{verbatim}
┌──────────────────────────────────────┐
│        Empleado (Abstract)           │
├──────────────────────────────────────┤
│ - id: long                           │
│ - name: String                       │
│ - position: String                   │
│ - email: String                      │
│ - phone: String                      │
├──────────────────────────────────────┤
│ + registrarAsistencia(): void        │
│ + getters/setters                    │
└──────────────┬───────────────────────┘
               │ heredan
      ┌────────┼─────────────┐
      │        │             │
      ▼        ▼             ▼
  ┌────────┐┌────────┐┌──────────────┐
  │Cocinero││ Mozo   ││PersonalLimpz.│
  └────────┘└────────┘└──────────────┘
\end{verbatim}

\subsection{Código: Clase Base Abstracta}

\inputminted[firstline=1, lastline=35]{java}{../../src/main/java/com/wallpawawqi/Class/Empleado.java}

\subsection{Beneficios}

\begin{itemize}
    \item \textbf{Reutilización de código}: Las propiedades comunes se definen una sola vez
    \item \textbf{Cohesión}: Todos los empleados comparten la misma interfaz base
    \item \textbf{Extensibilidad}: Nuevos tipos de empleados se agregan fácilmente
    \item \textbf{Polimorfismo}: Se pueden tratar todos los empleados uniformemente
\end{itemize}

\newpage

\section{2. Strategy Pattern}

\subsection{Descripción}

El patrón Strategy define una familia de algoritmos, encapsula cada uno y los hace intercambiables. Permite que el algoritmo varíe independientemente de los clientes que lo usan.

\subsection{Aplicación en WallpaWawqi}

Las subclases de \texttt{Empleado} representan diferentes \textbf{estrategias} de trabajo:

\begin{itemize}
    \item \textbf{Cocinero}: Estrategia de preparación de platillos
    \item \textbf{Mozo}: Estrategia de atención al cliente
    \item \textbf{PersonalDeLimpieza}: Estrategia de mantenimiento
\end{itemize}

Cada estrategia implementa sus propios métodos específicos que definen cómo ejecuta sus funciones.

\subsection{Diagrama Conceptual}

\begin{verbatim}
┌─────────────────────────────────────┐
│         Strategy (Empleado)         │
│      Interfaz: registrarAsistencia()│
└─────────────┬───────────────────────┘
              │
    ┌─────────┼──────────┐
    │         │          │
    ▼         ▼          ▼
┌─────────┐┌─────────┐┌──────────────┐
│ Cooking ││ Service ││Maintenance   │
│Strategy ││Strategy ││Strategy      │
└─────────┘└─────────┘└──────────────┘
\end{verbatim}

\subsection{Código: Estrategia 1 - Cocinero}

\inputminted[firstline=1, lastline=22]{java}{../../src/main/java/com/wallpawawqi/Class/Subclases/Cocinero.java}

\subsection{Código: Estrategia 2 - Mozo}

\inputminted[firstline=1, lastline=26]{java}{../../src/main/java/com/wallpawawqi/Class/Subclases/Mozo.java}

\subsection{Código: Estrategia 3 - PersonalDeLimpieza}

\inputminted[firstline=1, lastline=19]{java}{../../src/main/java/com/wallpawawqi/Class/Subclases/PersonalDeLimpieza.java}

\subsection{Ventajas}

\begin{itemize}
    \item \textbf{Flexibilidad}: Las estrategias pueden cambiar en tiempo de ejecución
    \item \textbf{Modularidad}: Cada estrategia es independiente y auto-contenida
    \item \textbf{Testabilidad}: Las estrategias se pueden probar aisladamente
    \item \textbf{Mantenibilidad}: Los cambios en una estrategia no afectan otras
\end{itemize}

\newpage

\section{3. Service Layer Pattern}

\subsection{Descripción}

El patrón Service Layer encapsula la lógica de negocio y acceso a datos en una capa separada. Proporciona una interfaz clara para operaciones específicas del dominio.

\subsection{Aplicación en WallpaWawqi}

La clase \texttt{JsonPersistenceService<T>} actúa como servicio reutilizable que:

\begin{itemize}
    \item Maneja la carga de datos desde JSON
    \item Persiste datos en archivos JSON
    \item Administra la generación de IDs
    \item Abstrae la complejidad de I/O de archivos
\end{itemize}

\subsection{Responsabilidades}

\begin{verbatim}
JsonPersistenceService
├── load()    → Carga datos del archivo JSON
├── save()    → Persiste datos en JSON
└── add()     → Agrega nuevo elemento y asigna ID
\end{verbatim}

\subsection{Código: Clase Servicio}

\inputminted[firstline=1, lastline=40]{java}{../../src/main/java/com/wallpawawqi/services/JsonPersistenceService.java}

\subsection{Beneficios}

\begin{itemize}
    \item \textbf{Separación de responsabilidades}: Lógica de datos separada del dominio
    \item \textbf{Reutilización}: El servicio funciona con cualquier tipo de datos
    \item \textbf{Testabilidad}: El servicio puede mockarse fácilmente
    \item \textbf{Mantenibilidad}: Cambios en la persistencia no afectan el dominio
\end{itemize}

\newpage

\section{4. Generic Types Pattern}

\subsection{Descripción}

Los genéricos de Java permiten escribir código parametrizado por tipos, mejorando la reutilización y seguridad de tipos sin sacrificar la flexibilidad.

\subsection{Aplicación en WallpaWawqi}

\texttt{JsonPersistenceService<T>} utiliza un parámetro de tipo genérico \texttt{T} para:

\begin{itemize}
    \item Trabajar con cualquier clase sin duplicar código
    \item Mantener seguridad de tipos en tiempo de compilación
    \item Evitar casting explícito innecesario
\end{itemize}

\subsection{Ejemplo de Uso}

\begin{minted}{java}
// Para persistencia de Platillos
JsonPersistenceService<Platillo> servicePlatillos =
    new JsonPersistenceService<>();
List<Platillo> platillos = servicePlatillos.load(
    "platillos.json",
    Platillo[].class
);

// Para persistencia de Empleados
JsonPersistenceService<Empleado> serviceEmpleados =
    new JsonPersistenceService<>();
List<Empleado> empleados = serviceEmpleados.load(
    "empleados.json",
    Empleado[].class
);
\end{minted}

\subsection{Ventajas}

\begin{itemize}
    \item \textbf{DRY (Don't Repeat Yourself)}: Una sola implementación para múltiples tipos
    \item \textbf{Type Safety}: El compilador valida los tipos
    \item \textbf{Legibilidad}: El código es más claro y auto-documentado
    \item \textbf{Rendimiento}: Sin sobrecarga de casting dinámico
\end{itemize}

\newpage

\section{5. Functional Programming Pattern}

\subsection{Descripción}

La programación funcional trata las funciones como ciudadanos de primera clase, permitiendo pasar comportamiento como parámetros mediante interfaces funcionales o method references.

\subsection{Aplicación en WallpaWawqi}

El método \texttt{add()} de \texttt{JsonPersistenceService} utiliza este patrón:

\begin{minted}{java}
public long add(
    String filePath,
    List<T> data,
    T nuevo,
    IdSetter<T> idSetter,    // Comportamiento: set ID
    IdGetter<T> idGetter     // Comportamiento: get ID
) throws IOException {
    // Calcula nuevo ID
    long nuevoId = data.stream()
        .mapToLong(idGetter::getId)    // Method reference
        .max()
        .orElse(0) + 1;

    idSetter.setId(nuevo, nuevoId);
    data.add(nuevo);
    save(filePath, data);

    return nuevoId;
}
\end{minted}

\subsection{Interfaces Funcionales Definidas}

\inputminted[firstline=35, lastline=40]{java}{../../src/main/java/com/wallpawawqi/services/JsonPersistenceService.java}

\subsection{Uso en la Práctica}

En el controlador HTTP se utiliza así:

\begin{minted}{java}
long nuevoId = service.add(
    DATA_FILE,
    platillos,
    nuevoPlatillo,
    Platillo::setId,    // Method reference a setter
    Platillo::getId     // Method reference a getter
);
\end{minted}

\subsection{Beneficios}

\begin{itemize}
    \item \textbf{Flexibilidad}: El comportamiento se inyecta como parámetro
    \item \textbf{Reutilización}: La lógica de \texttt{add()} es agnóstica del tipo
    \item \textbf{Concisión}: Los method references hacen el código más legible
    \item \textbf{Expresividad}: El código describe \textit{qué} hacer, no \textit{cómo}
\end{itemize}

\newpage

\section{6. Dependency Injection (Implícito)}

\subsection{Descripción}

Aunque no es explícito mediante frameworks como Spring, el código utiliza patrones que simulan inyección de dependencias.

\subsection{Aplicación}

En \texttt{WallpaWawqi.main()}:

\begin{minted}{java}
Gson gson = new Gson();  // Inyección manual
JsonPersistenceService<Platillo> service =
    new JsonPersistenceService<>();  // Inyección manual

server.createContext("/platillos", exchange -> {
    // Uso de dependencias
    List<Platillo> platillos = service.load(...);
    String json = gson.toJson(platillos);
});
\end{minted}

Este patrón permite:
\begin{itemize}
    \item Desacoplamiento de dependencias
    \item Facilidad de testing (reemplazar con mocks)
    \item Flexibilidad en la configuración
\end{itemize}

\newpage

\chapter{Análisis Integrado}

\section{Interacción de Patrones}

Los patrones se integran de forma coherente:

\begin{enumerate}
    \item \textbf{Template Method + Strategy}: Define roles y comportamientos
    \item \textbf{Service Layer}: Abstrae la persistencia
    \item \textbf{Generics}: Hace el servicio reutilizable
    \item \textbf{Functional Programming}: Inyecta comportamientos específicos
\end{enumerate}

\section{Flujo de Datos}

\begin{verbatim}
HTTP Request
    ↓
WallpaWawqi (Main)
    ↓
JsonPersistenceService<T>
    ↓
Gson (JSON Serialization)
    ↓
File System (JSON)
\end{verbatim}

\section{Ejemplo Completo: Agregar un Platillo}

\inputminted[firstline=45, lastline=65]{java}{../../src/main/java/com/wallpawawqi/WallpaWawqi.java}

Este código demuestra:
\begin{itemize}
    \item \textbf{Generics}: \texttt{JsonPersistenceService<Platillo>}
    \item \textbf{Service Layer}: Delegación al servicio
    \item \textbf{Functional Programming}: \texttt{Platillo::setId}
    \item \textbf{Dependency Injection}: Instancias inyectadas
\end{itemize}

\newpage

\chapter{Evaluación y Recomendaciones}

\section{Fortalezas}

\begin{enumerate}
    \item \textbf{Arquitectura Modular}: Separación clara de responsabilidades
    \item \textbf{Reutilización de Código}: Service genérico evita duplicación
    \item \textbf{Extensibilidad}: Fácil agregar nuevos tipos de empleados
    \item \textbf{Buenas Prácticas}: Uso correcto de patrones GOF
    \item \textbf{Type Safety}: Uso de genéricos y tipos fuerte
\end{enumerate}

\section{Áreas de Mejora}

\subsection{1. Implementación de Métodos Plantilla}

\textbf{Recomendación}: Los métodos plantilla en \texttt{Empleado} tienen solo TODOs.

\begin{minted}{java}
// Actual
public void registrarAsistencia() {
    // TODO registrar asistencia del empleado
}

// Propuesto
public void registrarAsistencia() {
    System.out.println("Asistencia registrada para: " + name);
    // Persistir en base de datos
}
\end{minted}

\subsection{2. Persistencia de Empleados}

\textbf{Recomendación}: Extender \texttt{JsonPersistenceService} para persistir empleados.

\begin{minted}{java}
// Nuevo endpoint HTTP
server.createContext("/empleados", exchange -> {
    if ("GET".equals(exchange.getRequestMethod())) {
        JsonPersistenceService<Empleado> service =
            new JsonPersistenceService<>();
        List<Empleado> empleados = service.load(
            "src/main/resources/empleados.json",
            Empleado[].class
        );
        // Responder con JSON
    }
});
\end{minted}

\subsection{3. Patrón Data Access Object (DAO)}

\textbf{Recomendación}: Crear una interfaz DAO para abstraer mejor la persistencia.

\begin{minted}{java}
public interface Dao<T> {
    List<T> findAll() throws IOException;
    Optional<T> findById(long id) throws IOException;
    void save(T entity) throws IOException;
    void delete(long id) throws IOException;
}

public class JsonDao<T> implements Dao<T> {
    private JsonPersistenceService<T> service;
    private String filePath;

    @Override
    public List<T> findAll() throws IOException {
        return service.load(filePath, getClazz());
    }
}
\end{minted}

\subsection{4. Manejo de Errores}

\textbf{Recomendación}: Implementar excepciones personalizadas.

\begin{minted}{java}
public class DataPersistenceException extends Exception {
    public DataPersistenceException(String msg, Throwable cause) {
        super(msg, cause);
    }
}

public class InvalidEmployeeException extends Exception {
    public InvalidEmployeeException(String msg) {
        super(msg);
    }
}
\end{minted}

\subsection{5. Validación de Datos}

\textbf{Recomendación}: Agregar validación en el servicio.

\begin{minted}{java}
public long add(String filePath, List<T> data, T nuevo,
                IdSetter<T> idSetter, IdGetter<T> idGetter)
        throws IOException, InvalidDataException {
    if (nuevo == null) {
        throw new InvalidDataException("El objeto no puede ser nulo");
    }
    // ... resto del método
}
\end{minted}

\newpage

\chapter{Conclusiones}

\section{Resumen de Patrones}

El backend de WallpaWawqi implementa correctamente una combinación coherente de patrones de diseño:

\begin{table}[h]
\centering
\begin{tabular}{|l|l|l|}
\hline
\textbf{Patrón} & \textbf{Componente} & \textbf{Propósito} \\
\hline
Template Method & Empleado (abstract) & Estructura común \\
\hline
Strategy & Cocinero, Mozo, PersonalLimpieza & Comportamientos específicos \\
\hline
Service Layer & JsonPersistenceService & Abstracción de datos \\
\hline
Generic Types & JsonPersistenceService<T> & Reutilización sin duplicación \\
\hline
Functional & IdSetter/IdGetter & Inyección de comportamiento \\
\hline
Dependency Injection & Manual (main) & Desacoplamiento \\
\hline
\end{tabular}
\end{table}

\section{Eficacia Arquitectónica}

La arquitectura es \textbf{efectiva} porque:

\begin{itemize}
    \item Cumple con el principio de Single Responsibility
    \item Facilita la extensión sin modificar código existente (Open/Closed)
    \item Permite sustituir implementaciones (Dependency Inversion)
    \item Evita duplicación mediante genéricos
    \item Mantiene el código limpio y comprensible
\end{itemize}

\section{Aplicabilidad a Requisitos}

Los patrones seleccionados son \textbf{apropiados} para:

\begin{enumerate}
    \item \textbf{Gestión de múltiples roles}: Strategy facilita roles diferentes
    \item \textbf{Persistencia flexible}: Genéricos permiten múltiples tipos
    \item \textbf{API REST}: Service Layer expone lógica clara
    \item \textbf{Extensibilidad}: Nuevos tipos se agregan sin refactoring mayor
\end{enumerate}

\section{Recomendación Final}

Se recomienda mantener esta arquitectura de patrones como base, implementando las mejoras sugeridas en el Capítulo 5, particularmente:

\begin{itemize}
    \item Patrón DAO para persistencia más robusta
    \item Validación y manejo de excepciones mejorado
    \item Completar la implementación de métodos plantilla
    \item Extensión de servicios para persistencia de Empleados
\end{itemize}

Con estos ajustes, el sistema estará mejor preparado para evolución futura y mantenimiento a largo plazo.

\newpage

\appendix

\chapter{Referencias de Patrones}

\section{Patrones de Diseño GOF}

\begin{itemize}
    \item \textbf{Creational Patterns}: Abstract Factory, Builder, Factory Method, Prototype, Singleton
    \item \textbf{Structural Patterns}: Adapter, Bridge, Composite, Decorator, Facade, Flyweight, Proxy
    \item \textbf{Behavioral Patterns}: Chain of Responsibility, Command, Interpreter, Iterator, Mediator, Memento, Observer, State, Strategy, Template Method, Visitor
\end{itemize}

\section{Libros Recomendados}

\begin{itemize}
    \item Design Patterns: Elements of Reusable Object-Oriented Software (Gang of Four)
    \item Head First Design Patterns (Freeman \& Freeman)
    \item Effective Java (Joshua Bloch)
    \item Clean Code (Robert C. Martin)
\end{itemize}

\section{Recursos Online}

\begin{itemize}
    \item Refactoring.guru - Comprehensive Design Patterns
    \item Oracle Java Tutorials - Generics
    \item Java Functional Interfaces Documentation
\end{itemize}

\end{document}
